\section{Introduction}
\label{sec:intro}

Deterioration events in hospitalised patients unfold over hours, and the physiological trajectory that precedes them is recorded long before the event itself is recognised. This temporal structure has motivated a substantial literature on data-driven early warning, ranging from static severity scores to sequential models that forecast deterioration trajectories from electronic health records \citep{xiumin2026machine} and compact, efficiency-aware networks that predict sepsis onset from structured time series \citep{taparugssanagorn2026lightweight}. Reported discrimination is high and improving, with in-hospital mortality and onset prediction now routinely exceeding the performance of classical scoring instruments \citep{yong2024deep}. The implicit promise of this line of work is that a sufficiently accurate model, once installed at the bedside, will shorten time to intervention.

Discriminative prediction and operational alerting address different questions. A score-to-alert conversion requires a decision rule, whereas threshold-independent ranking metrics do not specify how frequently that rule emits alarms. Reviews of intensive-care adverse-event prediction discuss incomplete reporting of alert burden, time-to-alert and deployment calibration \citep{wang2026deep}. Reviews of multivariate clinical time-series learning also motivate evaluation beyond architectural comparisons \citep{hernandezdiaz2026deep,swinckels2024deep}. These concerns support reporting the alarm policy and its realised load explicitly, rather than treating an area-under-the-curve comparison as a specification of bedside alerting behaviour.

Recent research improves the backbone that produces the risk score. For irregularly sampled and heterogeneously missing observations, approaches include reliability-aware temporal aggregation \citep{kim2026retamamba,shen2026enhancing}, graph formulations over irregular observations \citep{yalavarthi2024grafiti}, and masked-attention models for deterioration prediction \citep{mun2026made}. Transfer-oriented variants align feature spaces across cohorts \citep{komolafe2025early}, while dynamic temporal graphs incorporate clinical calculators into the prediction pathway \citep{yin2024sepsiscalc}. These methods motivate prospective backbone comparisons, but their adapters are not included in the current reference implementation. Our implemented contribution instead concerns the causal mapping from fixed scores to stateful alarms and the risk associated with selecting that mapping.

A separate literature attaches reliability estimates to clinical predictions. Selective prediction and calibration have been analysed for multimodal clinical classification \citep{lpez2026empirical}, reliability estimation has been extended beyond output-space calibration for time-series classifiers \citep{cenacchi2026beyond}, and conformal methods have been used for physiological regression uncertainty \citep{strickland2025degrees}. Our distinction is the unit and estimand of alarm-policy calibration: refractory state makes successive emitted alarms dependent, so those hours cannot simply be counted as independent patient samples. This does not imply that every conformal treatment of time series is invalid or that a stateful policy cannot be calibrated at the trajectory level. The reference construction uses independent patient blocks and fixed finite-grid upper-confidence bounds, a stronger sampling assumption than exchangeability alone.

This draft makes the alarm budget an explicit constraint and calibrates the policy that generates alerts. The implementation separates causal score estimation from full-trajectory policy replay. Its finite-sample bound concerns the population mean of patient-level rates, not a hard maximum for each patient or ward. Real-cohort comparisons remain prospective, and the current contributions are as follows.
\begin{itemize}
\item We formulate early warning as a stateful alarm policy in which the refractory window and escalation tier are control variables, so that alarm volume becomes a design parameter of the monitoring system rather than an emergent property of a threshold.
\item We derive a patient-blocked finite-grid upper-confidence bound on the population alarm rate under independent sampling, and distinguish it from a deliberately invalid timestep effective-sample-size shortcut.
\item We provide causal reference learners and official-format data readers, and specify a prospective evaluation protocol under shared requested budgets, with cross-hospital transfer and observation sparsification as planned extensions.
\end{itemize}
The remainder of the paper reviews related work (\cref{sec:related}), develops the policy formulation and the calibration procedure (\cref{sec:method}), and specifies the datasets, baselines, metrics, and experimental protocol (\cref{sec:setup}).
